\section{Design}\label{sec:design}

\begin{figure*}[t]
  \centering
  \resizebox{\textwidth}{!}{\input{figures/sequence}}
  \Description{A message sequence between six parties: the wallet in the user's browser, the user's mail provider, the sealed vault, the Perch recovery controller, the per-account email guardian, and the DKIM key registry with the proof verifier. In enrollment, the wallet draws an account code, the user sends an enrollment email that the provider signs and the vault seals back to the browser, the wallet checks it and computes the commitment and the pin, deploys the guardian, and applies a document listing it. In recovery, an attempt is opened, the controller returns the statement and its digest, the wallet draws rho and computes the blinded value T, the user sends an email whose Subject ends with T, the provider signs it and the vault seals it back, the browser proves the circuit's obligations, the wallet submits the approval, the controller calls the guardian's check, which asks the registry and the verifier, and records one vote. Other guardians reach the quorum, the attempt is authorized, the alert service sends an email and a Web Push with a code, and after the delay, unless the owner cancels, the recovery completes.}
  \caption{Enrollment and recovery. Dashed arrows are answers. The proof is
    made in the browser; the vault, an attested enclave, is the only party
    of the operator's that sees an email or an address.}
  \label{fig:sequence}
\end{figure*}

Figure~\ref{fig:sequence} shows the two flows. Each account that opts in gets
its own immutable email-guardian contract, listed as one guardian in its
Perch recovery member and counted toward the quorum like any other. The
guardian fixes at construction the commitment $\Cm$ to the enrolled address,
the account, its Perch controller, the DKIM key registry it trusts, an
optional key pin, the actions it may approve, and the proof verifier.

\subsection{Statement binding}\label{sec:design-statement}
The email factor approves exactly Perch's recovery statement and nothing
else; there is no email-specific statement. The guardian authorizes
$\digest = \Sha(S)$ for the statement encoding $S$, the circuit binds
$\digest$, and the guardian also receives $S$ itself. The four items a
recovery must bind are fields of $S$:
\begin{description}
  \item[Account] the \texttt{account} contract id, which the guardian
    requires to be its own.
  \item[New key] for lost-key and compromise recovery, the hash of the
    replacement credentials and of the derived target document; for a
    reconfiguration the new configuration hash, for an upgrade the Wasm hash.
  \item[Nonce] the attempt id, which Perch never reuses for an account
    (request id for upgrades), and the configuration epoch, which every
    configuration change and completion bumps.
  \item[Expiry] \texttt{valid\_until\_ledger}, a ledger sequence number;
    Perch's statement carries no wall-clock time, so the DKIM \texttt{t=}
    tag is not used.
\end{description}
The action is inside $\digest$, so one Subject format serves every action,
and the guardian checks the action against its allowed set. The circuit takes
the digest as two 128-bit limbs $\digest_{hi}, \digest_{lo}$, exactly as
Perch's ZK projection splits it, so both are always canonical field elements.

\subsection{The recovery email and its blinding}\label{sec:design-email}
The user sends the recovery email directly from the enrolled inbox to
\textsf{RECOVERY\_INBOX}, one fixed address per deployment; nothing about
the user is needed to start recovery. The wallet opens the user's mail client
with the recipient and Subject filled in. The Subject is exactly
\begin{gather*}
  \texttt{APPROVE NIDO ACCOUNT RECOVERY } \mathrm{hex}(T),\\
  T = \blindf(\digest, \blind) = \Poseidon(\mathsf{DOM\_BLIND}, \digest_{hi}, \digest_{lo}, \blind),
\end{gather*}
where $\blind$ is 31 fresh random bytes that never leave the browser and
$\mathrm{hex}$ gives 64 lowercase hexadecimal characters. Every statement and
its digest are public, so a Subject carrying $\digest$ in the clear would let
the provider match the email to the account and attempt; $T$ hides
$\digest$ from anyone without $\blind$ when Poseidon2 is modelled as a random
function\claim{blinded-subject-hiding}. Nothing may precede or follow the
value: a \texttt{Re:} prefix, an RFC~2047 encoded word~\cite{rfc2047} or
trailing text makes the email unprovable, while a provider that folds the
long line at a space is harmless, because relaxed canonicalization restores
the exact bytes\claim{subject-blinding}.

The proof binds neither the recipient, nor the body, nor \texttt{t=}: not
every launch provider signs \texttt{To:} or sets \texttt{t=}, an email to any
recipient is equally good evidence of inbox control, and freshness comes from
the statement's deadline and single-use attempt id. The signature still
covers the whole signed header block.

The signed octets reach the browser through the sealed vault, which runs
the recovery address's mail exchanger: the browser first registers $T$ with
a fresh device key, sealed to the vault with HPKE~\cite{rfc9180}; the first
registration of a value holds it; and the vault seals each matching message
to that key and keeps nothing\claim{sealed-inbound}. Alternatively the user
imports the sent message as an \texttt{.eml} file.

\subsection{Enrollment and the commitment}\label{sec:design-enroll}
While the owner still holds the passkey, the wallet draws $\code$, 16 bytes
from a secure generator, never reused and never sent to Nido. The user sends
an enrollment email whose Subject is a distinct phrase followed by
$\mathrm{hex}(E)$, where $E$ hashes the account's contract id under its own
tag with a fresh $\blind_e$, and the wallet reads the exact
\texttt{From} address $\addr$ from the DKIM-signed copy, so the user never
types it and the provider writes the same bytes on every later email. The
commitment is
\begin{align*}
  \Cm = \commitf(\addr, \code) = \Poseidon(&\mathsf{DOM\_COMMIT}, |\addr|,\\
    &\mathrm{pack}(\addr), \code),
\end{align*}
where $\mathrm{pack}$ cuts the address, right-padded to 124 bytes, into four
31-byte field elements and each $\mathsf{DOM}$ is a distinct domain tag; the domain of the address is inside $\Cm$, so the
enrolled domain is the committed address's. The wallet checks the enrollment
email under the recovery circuit's own rules before deploying a guardian, so
an inbox whose recovery email could never be proven is refused at
enrollment\claim{enrollment-email-check}; an enrollment Subject can never
satisfy the recovery circuit\claim{enrollment-subject-distinct}. Nothing on
chain verifies the enrollment email: a wrong $\Cm$ disables only its own
owner's email factor.

The wallet then deploys the guardian with $\Cm$ and a key pin
(\S\ref{sec:design-registry}), reads back its configuration and code hash,
and applies a document that lists it. Listing is an ordinary Perch
reconfiguration, so under \texttt{Protected} it needs the current recovery
condition too; since the guardian is immutable, changing the address, the
code, the pin or the allowed actions means deploying a new guardian through
the same gate.

Recovery needs $\code$, so the wallet offers two backups: a block in the
enrollment email's body, kept in the Sent folder, encrypting $\code$ under an
Argon2id~\cite{rfc9106} key from a wallet-generated five-word passphrase; and
a recovery kit\claim{account-code-backup}.

\subsection{The circuit relation}\label{sec:design-circuit}
The recovery circuit (Fig.~\ref{fig:relation}) has four public inputs, in
this order: $\Cm$, $\digest_{hi}$, $\digest_{lo}$, and $\mathit{root}$, the
root of the key tree the signing key is proven against. The guardian builds
all four itself. The private inputs are the signed header block $H$ (at most
1536 bytes), the RSA signature $s$ and modulus $n$, $\blind$, $\code$, and the
key's leaf index and 16 Merkle siblings. Positions inside $H$ are hints that
the circuit computes in unconstrained code and then checks. There is no
action input, timestamp or nullifier: the action and the expiry are in the
digest, and Perch's controller, not a nullifier, stops double counting.

\begin{figure}[t]
  \centering
  \resizebox{\columnwidth}{!}{\input{figures/circuit-relation}}
  \Description{Three columns. Private inputs on the left: the signed header block H of at most 1536 bytes, the RSA signature and modulus, the account code, rho, and the leaf index with 16 siblings. Obligations C1 to C6 in the middle. Public inputs on the right: the commitment, the two digest limbs, and the key-tree root. Arrows show that the header feeds C1 to C5, the signature and modulus feed C1 and C6, the account code feeds C3, rho feeds C5, the Merkle path feeds C6, and C3, C5 and C6 open the commitment, the digest limbs and the root.}
  \caption{The recovery circuit's relation. Domain tags and lengths inside
    the Poseidon2 hashes are omitted.}
  \label{fig:relation}
\end{figure}

The circuit enforces six obligations:
\begin{description}
  \item[C1] $s<n$, $n$ has exactly 2048 bits, and $s^{65537} \bmod n$ is the
    EMSA-PKCS1-v1\_5 encoding of $\Sha(H)$.
  \item[C2] $H$ is a well-formed signed block: header fields in relaxed
    canonical form, each ended by CRLF, then the signature's own field with
    an empty \texttt{b=} and no trailing CRLF, carrying
    \texttt{v=1}, \texttt{a=rsa-sha256}, relaxed header canonicalization,
    an \texttt{h=} naming \texttt{from} and \texttt{subject}, and a
    \texttt{d=}; no byte past $|H|$ is read.
  \item[C3] $H$ holds exactly one \texttt{from} field; its value parses to an
    address $\addr$ that, with $\code$, opens $\Cm$.
  \item[C4] The lowercased domain of $\addr$ equals the lowercased
    \texttt{d=}, which is the domain of the registry leaf.
  \item[C5] $H$ holds exactly one \texttt{subject} field, whose value is the
    recovery phrase followed by $\mathrm{hex}(T)$ for the public digest limbs
    and $\blind$.
  \item[C6] The leaf $\leaff(\mathit{dom}, \kh(n))$ has a depth-16 Merkle
    path to $\mathit{root}$, where $\kh(n) = \Poseidon(\mathsf{DOM\_KEY},
    n_0,\ldots,n_{17})$ over the 18 120-bit limbs of $n$, and the leaf hashes
    $\mathsf{DOM\_LEAF}$, the domain's length, the domain packed into two
    field elements, and $\kh(n)$.
\end{description}
The circuit takes the canonicalized bytes the signature covers and never
canonicalizes; a field outside $H$ has no effect. The \texttt{From} value
must be a bare address or a printable display name followed by exactly one
angle-bracketed address that ends the field; quoted local parts, comments,
groups and address lists are refused, the address has at most 124 bytes,
and it is used exactly as signed, with no case folding or tag
removal\claim{circuit-from-grammar}. The circuit refuses simple header
canonicalization, a missing or duplicated \texttt{From} or \texttt{Subject},
a misaligned \texttt{d=}, and a signature not below its
modulus\claim{circuit-dkim-rules}. Changing any public input, or any private
value that links the email to one, makes the proof
fail\claim{circuit-binds-public-inputs}, and an honest recovery email from
every launch domain is accepted, including a 124-byte address with a 500-byte
display name\claim{circuit-accepts-launch-providers}.

\subsection{The email guardian}\label{sec:design-guardian}
The guardian's constructor refuses an account or controller that is not a
contract, the account as its own controller, a zero or non-canonical
commitment, a pin that has expired or lasts more than 365 days, an empty or
unknown action set, and a verifier whose reported circuit id differs from the
configured one\claim{guardian-construction}. Perch's controller calls
\texttt{guardian.require\_auth\_for\_args((d,))}, and the host then runs the
guardian's \texttt{\_\_check\_auth} with the approval $(S, \mathit{root},
\pi)$ from the guardian's authorization entry. The guardian refuses at the
first failing check of Table~\ref{tab:guardian}, and builds the public
inputs only from its configuration, the authorized digest and the accepted
root\claim{guardian-checks}. Approvals by the email guardian count toward a
\texttt{Protected} reconfiguration through \texttt{approve\_change} like
any guardian's\claim{guardian-on-perch}.

\begin{table}[t]
  \caption{The guardian's checks, in order.}
  \label{tab:guardian}
  \small
  \begin{tabularx}{\columnwidth}{@{}r>{\raggedright\arraybackslash}X l@{}}
    \toprule
    \# & Check & Error \\
    \midrule
    1 & One contract call, on the controller, with one 32-byte argument $\digest$ & \texttt{UnexpectedContext} \\
    2 & $\Sha(S) = \digest$ & \texttt{DigestMismatch} \\
    3 & $S$ is a version-2 Perch statement of a known action & \texttt{MalformedStatement} \\
    4--6 & $S$ names this network, account and controller & \texttt{Wrong\ldots} \\
    7 & The action is allowed & \texttt{ActionNotAllowed} \\
    8 & The call is that action's entry point & \texttt{WrongEntryPoint} \\
    9 & Ledger $\le$ \texttt{valid\_until\_ledger} & \texttt{StatementExpired} \\
    10 & $\mathit{root}$ is canonical & \texttt{NonCanonicalRoot} \\
    11 & $\mathit{root}$ is the unexpired pin, or the registry accepts it & \texttt{UnknownRoot} \\
    12 & $\pi$ verifies for $(\Cm, \digest_{hi}, \digest_{lo}, \mathit{root})$ & \texttt{ProofRejected} \\
    \bottomrule
  \end{tabularx}
\end{table}

\paragraph{One guardian per account.} Perch binds guardian addresses into
\texttt{config\_hash}, so everything fixed at the guardian's address is
covered by Perch's reconfiguration gate. A guardian shared by many accounts
would need a table from account to commitment, and whoever could change that
table could change an account's email factor without passing the gate. And
\texttt{require\_auth\_for\_args} hands the guardian only $\digest$: it
learns which account $\digest$ concerns only from $S$ (check~2), and approves
only for its own (check~5).

\paragraph{What the guardian does not check.} It ignores the host's
signature payload (the entry's nonce, expiration and invocation hash): the
proof binds the statement, not an authorization entry, so anyone can re-wrap
an approval, and the re-wrapped copy can still authorize only the same
digest, which the controller counts once\claim{guardian-payload-irrelevant}.
Quorum, delay, cancellation and epoch are the controller's; the guardian
could not read them anyway, since Soroban refuses re-entry into the
controller that is calling it. The guardian calls only the registry and the
verifier, which embeds Perch's UltraHonk verifier with the circuit's
verification key\claim{verifier-contract}.

\subsection{The DKIM key registry and the pin}\label{sec:design-registry}
Keys are identified by $\kh(n)$, never by selector, because a provider may
reuse a selector for a new key. The registry holds a depth-16 Merkle tree of
leaves $\leaff(\mathit{domain}, \kh)$ hashed exactly as C6 hashes them, with the
host's Poseidon2, and answers one question for the guardian: is this root
accepted now\claim{registry-hash-parity}? Its rules are:
\begin{itemize}
  \item \emph{Activation.} An addition, proposed and approved by the
    multisig, executes 72 hours after the later of its approval and the
    ledger since which a threshold of attesters' latest observations see the
    key live in DNS, or 7 days after its approval, whichever comes
    first\claim{registry-timelocks}.
  \item \emph{Grace.} A root superseded at ledger $L$ stays acceptable
    through $L + \mathit{grace}$, so a proof against the previous root still
    verifies and a key the provider has just retired keeps working briefly.
  \item \emph{Revocation.} An emergency threshold of members revokes a key
    for compromise; from that ledger every root containing it is refused,
    grace included, a pending addition of it stops, and it can return only
    through a new full-threshold addition on the 7-day
    path\claim{registry-revocation}.
  \item \emph{Scope.} Leaves exist only for allowlisted domains and RSA keys
    of exactly 2048 bits; the registry computes key hashes itself; the empty
    tree's root is never accepted\claim{registry-governance}.
\end{itemize}
An attester resolves \texttt{selector.\_domainkey.domain}, signs an exact
encoding of what it saw with Ed25519~\cite{rfc8032}, bound to the registry's
contract id, and anyone may submit the signature; each attester's latest
observation counts\claim{registry-attestations}. The registry has no admin,
setter or upgrade path\claim{registry-immutable}: changing a parameter
means deploying a new registry, which a guardian follows only by being
replaced through Perch's gate. Its proposed governance is five members with
threshold three, an emergency threshold of two, three attesters with
threshold two, and a 24-hour grace period\claim{registry-defaults}.

\paragraph{The per-account pin.} A guardian may carry one pin, a key tree
fixed at construction, with an expiry of at most 365 days. Its leaf 0 is
$\Poseidon(\mathsf{DOM\_PIN\_SALT}, \code)$ and its other leaves are
registry-format leaves of the pinned keys, so the pin's root reveals nothing
about which provider's keys it holds without $\code$. The guardian accepts a
key root if it is the unexpired pin's or the registry accepts it, so the pin
carries a new account through the window while its provider's key waits in
the registry's timelock, and keeps the enrolled key usable if the registry
drops it\claim{key-pin}. The pin is the one path the registry cannot close: a
revoked key stays usable through a pin that holds it until the pin expires,
and forging an approval with it still needs the victim's address and
$\code$. The wallet deploys guardians with a 30-day
pin\claim{wallet-pin-lifetime}.

\paragraph{Finding the key.} Asking for one provider's key reveals the
provider, so the wallet reads a static list of every registry key, the same
for every user, and falls back to DNS over HTTPS only after telling the
user; either way a key counts only through the pin or an accepted
root\claim{key-lookup}.

\subsection{Alerts and the sealed vault}\label{sec:design-alerts}
Alerts are driven by the controller's own events; nothing is written on
chain for them. The alert service's host watches the controller by contract
id and hands each event to the vault, which finds the account's record under
$\mathrm{HMAC}(k_\mathit{alert}, \mathit{account})$ with a key that never
leaves it. Authorization, completion and configuration changes alert every
channel; opening an attempt, which anyone can do, at most one email a day.
An alert has generic text and no link, and carries a six-digit code keyed by
a key derived with HKDF~\cite{rfc5869} from $\code$, which the wallet shows
for the attempt it knows\claim{alert-code}.

When email is a recovery factor, a second channel is required: Web Push to a
passkey device~\cite{rfc8030,rfc8291,rfc8292}. A contact is registered,
replaced or removed only with an assertion by one of the account's
administrative passkeys~\cite{webauthn3}; a replacement or removal takes
effect after 14 days, longer than the 7-day minimum delay, while the old
contact keeps receiving alerts and is told at once; and removal is refused
while a recovery is authorized\claim{alert-contacts}. Alert email goes
through an EU-owned sending service, with failover to a second
one\claim{alert-formats}.

The browser encrypts each address to the vault with HPKE after verifying the
vault's attestation document~\cite{aws-nitro-attestation}. The vault's master
key is split 2-of-3 with Shamir's scheme~\cite{shamir} across independent
custodians, who release a share only to an enclave image that an on-chain
image registry has activated after a 14-day public timelock, and stop once
the image is revoked\claim{vault-custody}. With every third party mocked and
the vault in its own process, no message, host record, log, RPC call,
custodian decision or push body contains the address: Nido's servers and
staff can't read your address\claim{servers-cannot-read-address}, under A8.
The sending service and the user's mail provider necessarily see it.

\subsection{Wallet policy}\label{sec:design-policy}
Perch's controller accepts any configuration its owner applies, so the
remaining rules bind Nido's wallet, which builds every recovery member
through one builder that tells the email guardian apart by a typed marker,
never by address:
\begin{enumerate}
  \item at most one email guardian per recovery
    member\claim{wallet-one-email-guardian};
  \item an email guardian only with a registered second alert
    channel\claim{wallet-second-alert-channel};
  \item if the email vote alone can authorize (guardian-only, quorum one),
    the \texttt{Loss} profile and a delay of at least seven days, 120{,}960
    ledgers\claim{wallet-email-plus-delay};
  \item every action allowed whenever the quorum needs the email vote, so
    that a \texttt{Protected} account can still reconfigure and
    upgrade\claim{wallet-email-allowed-actions}.
\end{enumerate}
OIDC sign-in with Google, Apple or Microsoft may later become a second way
to produce the email factor; it would never count as a factor independent of
the same provider's DKIM, since both rest on one provider.
